%  LaTeX support: latex@mdpi.com 
%  For support, please attach all files needed for compiling as well as the log file, and specify your operating system, LaTeX version, and LaTeX editor.

%=================================================================
\documentclass[jimaging,article,submit,pdftex,moreauthors]{Definitions/mdpi} 


\graphicspath{{figures/}} % path to folder with figures
\usepackage{caption}
\usepackage{subcaption}
%\usepackage[export]{adjustbox}
\usepackage{listings}
\usepackage{xcolor}
\usepackage{graphicx}

\definecolor{codegreen}{rgb}{0,0.6,0}
\definecolor{codegray}{rgb}{0.5,0.5,0.5}
\definecolor{codepurple}{rgb}{0.58,0,0.82}
%\definecolor{backcolour}{rgb}{0.95,0.95,0.92}
\definecolor{backcolour}{RGB}{255,250,240}

\lstdefinestyle{mystyle}{
    backgroundcolor=\color{backcolour},   
    commentstyle=\color{codegreen},
    keywordstyle=\color{magenta},
    numberstyle=\tiny\color{codegray},
    stringstyle=\color{codepurple},
    basicstyle=\ttfamily\footnotesize,
    breakatwhitespace=false,         
    breaklines=true,                 
    captionpos=t, 
    keepspaces=true,                 
    numbers=left,                    
    numbersep=5pt,                  
    showspaces=false,                
    showstringspaces=false,
    showtabs=false,                  
    tabsize=2
    }

\lstset{style=mystyle}
\captionsetup[lstlisting]{position=top, margin=0cm,labelfont={bf, small, stretch=1.17}, labelsep=colon, textfont={small, stretch=1.17}, aboveskip=6pt, singlelinecheck=off, justification=justified}
	
\usepackage{soul} % highlighting text

%--------------------
% Class Options:
%--------------------
%----------
% journal
%----------

%---------
% article
%---------

%=================================================================
% MDPI internal commands
\firstpage{1} 
\makeatletter 
\setcounter{page}{\@firstpage} 
\makeatother
\pubvolume{1}
\issuenum{1}
\articlenumber{0}
\pubyear{2022}
\copyrightyear{2022}
%\externaleditor{Academic Editor: Firstname Lastname}
\datereceived{} 
%\daterevised{} % Only for the journal Acoustics
\dateaccepted{} 
\datepublished{} 
%\datecorrected{} % Corrected papers include a "Corrected: XXX" date in the original paper.
%\dateretracted{} % Corrected papers include a "Retracted: XXX" date in the original paper.
\hreflink{https://doi.org/} % If needed use \linebreak
%\doinum{}
%------------------------------------------------------------------
% The following line should be uncommented if the LaTeX file is uploaded to arXiv.org
%\pdfoutput=1

%=================================================================
% Add packages and commands here. The following packages are loaded in our class file: fontenc, inputenc, calc, indentfirst, fancyhdr, graphicx, epstopdf, lastpage, ifthen, lineno, float, amsmath, setspace, enumitem, mathpazo, booktabs, titlesec, etoolbox, tabto, xcolor, soul, multirow, microtype, tikz, totcount, changepage, attrib, upgreek, cleveref, amsthm, hyphenat, natbib, hyperref, footmisc, url, geometry, newfloat, caption

%=================================================================
% Full title of the paper (Capitalized)
\Title{Python Library \emph{EarthPy} for Satellite Image Processing and Terrain Analysis of Côte d'Ivoire, West Africa}

% MDPI internal command: Title for citation in the left column
\TitleCitation{Python Library \emph{EarthPy} for Satellite Image Processing and Terrain Analysis of Côte d'Ivoire, West Africa}

% Author Orchid ID: enter ID or remove command
\newcommand{\orcidauthorA}{0000-0002-5759-1089} % Add \orcidA{} behind the author's name
%\newcommand{\orcidauthorB}{0000-0000-0000-000X} % Add \orcidB{} behind the author's name

% Authors, for the paper (add full first names)
\Author{Olivier Debeir $^{1}$ and Polina Lemenkova $^{1}$\orcidA{}}
% \dagger,\ddagger
%\longauthorlist{yes}

% MDPI internal command: Authors, for metadata in PDF
\AuthorNames{Olivier Debeir $^{1}$ and Polina Lemenkova $^{1}$}

% MDPI internal command: Authors, for citation in the left column
\AuthorCitation{Debeir, O.; Lemenkova, P.}
% If this is a Chicago style journal: Lastname, Firstname, Firstname Lastname, and Firstname Lastname.

% Affiliations / Addresses (Add [1] after \address if there is only one affiliation.)
\address{%
$^{1}$ \quad Université Libre de Bruxelles, École polytechnique de Bruxelles (EPB, Brussels Faculty of Engineering), Laboratory of Image Synthesis and Analysis, Building L, Campus de Solbosch, ULB -- LISA CP165/57, Avenue Franklin D. Roosevelt 50, B-1050, Brussels, Belgium; polina.lemenkova@ulb.be or pauline.lemenkova@gmail.com}

% Contact information of the corresponding author
\corres{Correspondence: polina.lemenkova@ulb.be; Tel.: +32471860459 (P.L.)}

% Abstract (Do not insert blank lines, i.e. \\) 
\abstract{A single paragraph of about 200 words maximum.}

% Keywords
\keyword{Image processing, remote sensing; programming; Python; satellite image; Digital Terrain Model; cartography; mapping; spatial analysis; environment} 

% The fields PACS, MSC, and JEL may be left empty or commented out if not applicable
%\PACS{J0101}
%\MSC{}
%\JEL{}

%%%%%%%%%%%%%%%%%%%%%%%%%%%%%%%%%%%%%%%%%%
\begin{document}

\section{Introduction}
\newpage
\section{Study Area}

\begin{figure}[H]\centering
	\includegraphics[width=12.0 cm]{Fig_01}
	\caption{Topographic map of Côte d'Ivoire. Mapping software: Generic Mapping Tools (GMT) scripting toolset. Data source: GEBCO/SRTM. Cartography: authors.
	\label{fig1}}
\end{figure}  

%%%%%%%%%%%%%%%%%%%%%%%%%%%%%%%%%%%%%%%%%%
\section{Materials and Methods}

\begin{lstlisting}[language=python, caption=EarthPython code (1) used for plotting DEM in Figure 2.]
python3.10
# Set home directory
import os
os.chdir('/Users/polinalemenkova/Documents/Python/Image_Processing')
os. getcwd()
# Set data source
dtm = "n07_w006_3arc_v2.tif"
# Load libraries
from glob import glob
import earthpy as et
import earthpy.spatial as es
import earthpy.plot as ep
import rasterio as rio
import matplotlib.pyplot as plt
import numpy as np
from matplotlib.colors import ListedColormap
import plotly.graph_objects as go
np.seterr(divide='ignore', invalid='ignore')
# Open DEM by Rasterio which reads GeoTIFF format and stores gridded raster datasets of digital terrain models (DTMs) and satellite images
with rio.open(dtm) as src:
    elevation = src.read(1)
# Plot the data
fig, ax = plt.subplots(figsize=(10, 10))
ep.plot_bands(
    elevation,
    ax=ax,
    cmap="terrain",
    title="Ditigal Terrain Model (DTM) of Kossou Lake, C\^ote d'Ivoire. \n Data source: SRTM3 (3 arc-second resolution, 90 m), NASA/NGA. \nEntity ID: SRTM3N07W006V2. Aquisition date: 2000-02-11. cpt: 'terrain'",
    figsize=(10, 6),
)
plt.show()
# Save figure with 300 dpi resolution
plt.savefig('Figure_2_DEM_Kossou_1609.jpg', dpi=300)
\end{lstlisting}

\begin{lstlisting}[language=python, caption=EarthPython code (2) used for plotting hillshade in Figure 3.]
python3.10
# Set home directory
import os
os.chdir('/Users/polinalemenkova/Documents/Python/Image_Processing')
os. getcwd()
# Set data source
dtm = "n07_w006_3arc_v2.tif"
# Load libraries
from glob import glob
import earthpy as et
import earthpy.spatial as es
import earthpy.plot as ep
import rasterio as rio
import matplotlib.pyplot as plt
import numpy as np
from matplotlib.colors import ListedColormap
import plotly.graph_objects as go
np.seterr(divide='ignore', invalid='ignore')
# Generate hillshade from DTM
hillshade = es.hillshade(elevation)
# Plot the data
fig, ax = plt.subplots(figsize=(10, 10))
ep.plot_bands(
    hillshade,
    ax=ax,
    cbar=True,
    cmap="cividis",
    title="Hillshade from DTM of Kossou Lake, C\^ote d'Ivoire. \n Data source: SRTM3 (3 arc-second resolution, 90 m), NASA/NGA. \nEntity ID: SRTM3N07W006V2. Aquisition date: 2000-02-11. cmap=cividis",
    figsize=(10, 10),
)
plt.show()
plt.savefig('Figure_3_Hillshade_Kossou_1609.png', dpi=300)
\end{lstlisting}

\begin{lstlisting}[language=python, caption=EarthPython code (3) and (4) used for plotting hillshade with changed azimuth and sun angle values for Figures 4 and 5.]
python3.10
# Set home directory
import os
os.chdir('/Users/polinalemenkova/Documents/Python/Image_Processing')
os. getcwd()
# Set data source
dtm = "n07_w006_3arc_v2.tif"
# Load libraries
from glob import glob
import earthpy as et
import earthpy.spatial as es
import earthpy.plot as ep
import rasterio as rio
import matplotlib.pyplot as plt
import numpy as np
from matplotlib.colors import ListedColormap
import plotly.graph_objects as go
np.seterr(divide='ignore', invalid='ignore')
# Open DEM by Rasterio which reads GeoTIFF format and stores gridded raster datasets of digital terrain models (DTMs) and satellite images
with rio.open(dtm) as src:
    elevation = src.read(1)
# Changing the azimuth of the sun of the hillshade layer to 210 degrees
hillshade_azimuth_230 = es.hillshade(elevation, azimuth=230)
# Plotting the hillshade layer with the modified azimuth
fig, ax = plt.subplots(figsize=(10, 10))
ep.plot_bands(
    hillshade_azimuth_230,
    ax=ax,
    cbar=True,
    cmap="plasma",
    title="Hillshade from DTM of Kossou Lake, C\^ote d'Ivoire. Azimuth of sun is adjusted to 230 \N{DEGREE SIGN}",
    figsize=(10, 10),
)
plt.show()
plt.savefig('Figure_4_Azimuth_hillshade_Kossou.jpg', dpi=300)
# Changing the angle altitude of sun
hillshade_angle_10 = es.hillshade(elevation, altitude=10)
# Plotting the hillshade layer with the modified angle altitude
fig, ax = plt.subplots(figsize=(10, 10))
ep.plot_bands(
    hillshade_angle_10,
    ax=ax,
    cbar=True,
    cmap="magma",
    title="Hillshade from DTM of Kossou Lake, C\^ote d'Ivoire. Angle altitude is defined as 10 \N{DEGREE SIGN}",
    figsize=(10, 10),
)
plt.show()
plt.savefig('Figure_5_Angle_hillshade_Kossou.jpg', dpi=300)

\end{lstlisting}

%%%%%%%%%%%%%%%%%%%%%%%%%%%%%%%%%%%%%%%%%%
\section{Results}

\begin{figure}[H]\centering
	\includegraphics[width=12.0 cm]{Fig_02}
	\caption{Ditigal Terrain Model (DTM) of Kossou Lake, Côte d'Ivoire. Data source: SRTM3 (3 arc-sec resolution, 90 m), NASA/NGA. Entity ID: SRTM3N07W006V2; tile coordinates: 5°00'00"-6°00'00"W, 7°00'00"-8°00'00"N. Aquisition: 2000-02-11. cpt: 'terrain'. Mapping: Python. Source: authors.\label{fig2}}
\end{figure}  

\begin{figure}[H]\centering
	\includegraphics[width=12.0 cm]{Fig_03}
	\caption{Hillshade from DTM of Kossou Lake, Côte d'Ivoire. Data source: SRTM3 (3 arc-second resolution, 90 m), NASA/NGA. Entity ID: SRTM3N07W006V2. Aquisition date: 2000-02-11. cmap=cividis. Mapping: Python. Source: authors.\label{fig3}}
\end{figure} 

%\newpage
\begin{figure}
     \centering
     \begin{subfigure}[b]{0.7\textwidth}
         \centering
         \includegraphics[width=10cm]{Fig_04}
         \caption{Azimuth of sun is adjusted to 230 in hillshade from DTM.}
         \label{fig:04a}
     \end{subfigure} \vfill \vspace{2mm}
     \begin{subfigure}[b]{0.7\textwidth}
         \centering
         \includegraphics[width=10cm]{Fig_05}
         \caption{Angle altitude is defined as 10 in hillshade from DTM.}
         \label{fig:04a}
     \end{subfigure}
        \caption{Modified azimuth and sun angle degrees in hillshade of DTM of Kossou Lake, Côte d'Ivoire. Mapping: Python. Source: authors.}
        \label{fig4}
\end{figure}

\begin{figure}[H]\centering
	\includegraphics[width=12.0 cm]{Fig_06}
	\caption{SRTM3 Void Filled DEM (NASA/NGA) overlayed on top of a hillshade in Kossou Lake, Côte d'Ivoire, West Africa. cmaps: 'turbo' and 'gist\_yarg'. Mapping: Python. Source: authors.
	\label{fig6}}
\end{figure} 

\begin{figure}
     \centering
     \begin{subfigure}[b]{0.6\textwidth}
         \centering
         \includegraphics[width=10cm]{Fig_07a.png}
         \caption{Stretch improvement: 'linear'.}
         \label{fig:04a}
     \end{subfigure}
     \vfill \vspace{2mm}
     \begin{subfigure}[b]{0.6\textwidth}
         \centering
         \includegraphics[width=10cm]{Fig_07b.png}
         \caption{Stretch improvement: 'histogram'.}
         \label{fig:04a}
     \end{subfigure}
        \caption{Natural color composites of bands 2-3-4 with applied stretch option on 'plotRGB()' which improves the rendering of the scene: Kossou Lake, Côte d'Ivoire. Mapping: R. Source: authors}
        \label{fig7}
\end{figure}

\begin{figure}
     \centering
     \begin{subfigure}[b]{0.7\textwidth}
         \centering
         \includegraphics[width=10cm]{Fig_08a}
         \caption{False color composite visualization of bands 3-4-5: near-infrared (B5) as red, red (B4) as green, and green (B3) as blue.}
         \label{fig:04a}
     \end{subfigure} 
     \vfill \vspace{2mm}
     \begin{subfigure}[b]{0.7\textwidth}
         \centering
         \includegraphics[width=10cm]{Fig_08b}
         \caption{False color composite visualization of bands 7-6-5}
         \label{fig:04a}
     \end{subfigure}
        \caption{Landsat 8 OLI colour composite image: Kossou Lake surroundings, Côte d'Ivoire. Mapping: RStudio. Source: authors.}
        \label{fig8}
\end{figure}

\begin{figure}
\makebox[\linewidth][c]{%
\begin{subfigure}[b]{.6\textwidth}
\centering
\includegraphics[width=0.95\textwidth]{Fig_09a}
\caption{Bands 5-7-3.}
\end{subfigure}%
\begin{subfigure}[b]{.6\textwidth}
\centering
\includegraphics[width=0.95\textwidth]{Fig_09b}
\caption{Bands 7-4-2}
\end{subfigure}%
}\\
 \vfill \vspace{2mm}
\makebox[\linewidth][c]{%
\begin{subfigure}[b]{.6\textwidth}
\centering
\includegraphics[width=0.95\textwidth]{Fig_10a}
\caption{Color Infrared bands 5-4-3.}
\end{subfigure}%
\begin{subfigure}[b]{.6\textwidth}
\centering
\includegraphics[width=0.95\textwidth]{Fig_10b}
\caption{Bands 6-5-4}
\end{subfigure}%
}
\caption{False color composite visualization of Landsat 8 OLI image of Kossou Lake surroundings, Côte d'Ivoire. Mapping: RStudio. Source: authors.}
\end{figure}

\newpage
\subsection{Subsection}
\subsubsection{Subsubsection}


%%%%%%%%%%%%%%%%%%%%%%%%%%%%%%%%%%%%%%%%%%
\section{Discussion}


%%%%%%%%%%%%%%%%%%%%%%%%%%%%%%%%%%%%%%%%%%
\section{Conclusions}

%%%%%%%%%%%%%%%%%%%%%%%%%%%%%%%%%%%%%%%%%%

\funding{This research received no external funding.}

\institutionalreview{Not applicable}

\informedconsent{Not applicable}

\dataavailability{Not applicable} 

\acknowledgments{We thank the two anonymous reviewers for comments and reading of this manuscript}

\conflictsofinterest{The authors declare no conflict of interest} 

\abbreviations{Abbreviations}{
The following abbreviations are used in this manuscript:\\

\noindent 
\begin{tabular}{@{}ll}
CPT & Colour Palette Table \\
DCW & Digital Chart of the World \\
GEBCO & General Bathymetric Chart of the Oceans \\
GDAL & Geospatial Data Abstraction Library \\
GIS & Geographic Information System \\
GMT & Generic Mapping Tools \\
\end{tabular}
}

\reftitle{References}

% Please provide either the correct journal abbreviation (e.g. according to the “List of Title Word Abbreviations” http://www.issn.org/services/online-services/access-to-the-ltwa/) or the full name of the journal.
% Citations and References in Supplementary files are permitted provided that they also appear in the reference list here. 

%=====================================
% References, variant A: external bibliography
%=====================================
\bibliography{IvoryCoast.bib}
\nocite{*}

%%%%%%%%%%%%%%%%%%%%%%%%%%%%%%%%%%%%%%%%%%
%\end{adjustwidth}
\end{document}

